\documentclass[11pt,a4paper]{article}
\usepackage[margin=25mm]{geometry}
\usepackage{amsmath,amssymb,amsthm,hyperref}
\newtheorem{lemma}{Lemma}\newtheorem{theorem}[lemma]{Theorem}
\newcommand{\gen}[1]{\langle#1\rangle}
\title{Finite marked subgroups by successive stabilizers}
\author{Human Compression project}\date{}
\begin{document}\maketitle
Let $G$ be a specified finite group, let $J$ be its nonidentity
involutions, and let $A$ act on $G$ by automorphisms. The action need
not be faithful. A string is an ordered tuple $t=(t_1,\ldots,t_m)$
with $[t_i,t_j]=1$ whenever $|i-j|>1$. Write $H(t)=\gen t$.
Its intersection property is
\[
 \gen{t_i:i\in I}\cap\gen{t_i:i\in K}
       =\gen{t_i:i\in I\cap K}\quad(I,K\subseteq\{1,\ldots,m\}).
\]
A string satisfying this property is a C-marking of its generated
subgroup. It need not generate the ambient $G$.

\begin{lemma}[An exact extension test]\label{ip}
Suppose $t$ has the intersection property, $x\in J\setminus H(t)$,
and $x$ commutes with $t_1,\ldots,t_{m-1}$. The extended string
$(t_1,\ldots,t_m,x)$ has the intersection property if and only if
its suffix $(t_2,\ldots,t_m,x)$ has it and
\[
 H(t)\cap\gen{t_2,\ldots,t_m,x}=\gen{t_2,\ldots,t_m}.
\]
This gives a recursive exact test, ending at length at most two.
\end{lemma}
\begin{proof}
We retain the standard end-intersection criterion for string
markings: two end deletions with the intersection property, and
their correct common-middle intersection, imply the property for
the full string (McMullen--Schulte, Proposition 2E16).
The first end deletion here is the assumed old marking. In the
recursive suffix test its old prefix is a submarking of $t$, so
again has the property. Two distinct nonidentity involutions have
order-two cyclic subgroups with trivial intersection, giving the
base case. The condition $x\notin H(t)$ ensures distinctness in
every suffix. Induction now proves sufficiency. Necessity follows
by applying the full intersection property to the displayed
subsets and to every suffix.
\end{proof}

Let $A_t$ be the subgroup of $A$ fixing every entry of $t$, and put
\[
 P(t)=J\cap\bigcap_{i=1}^{m-1}C_G(t_i).
\]
The empty intersection is all of $J$. Starting at the empty tuple,
make the following recursion through a specified depth $d$.
At $t$, form $P(t)\setminus H(t)$, take one representative from
each $A_t$-orbit, and append each representative which passes
Lemma~\ref{ip}. Update $H(t)$ and the pointwise stabilizer.
Record a visited tuple whenever it has the desired invariant
property, and continue until depth $d$ or $H(t)=G$.

\begin{theorem}\label{coverage}
Before any optional pruning, this recursion visits precisely one
representative of every $A$-orbit of C-markings in $G$ of length
at most $d$. Thus it classifies any $A$-invariant selected family
among these markings. No extension of a tuple generating $G$ is
lost by stopping there.
\end{theorem}
\begin{proof}
The candidate set is invariant under $A_t$, since its conditions
are generation, membership and commutation with entries fixed
by $A_t$. Every C-marking is independent: if a mark lay in the
subgroup on the other marks, intersecting that subgroup with
the cyclic subgroup of the mark would contradict the property.
In particular each final mark lies outside its previous prefix
subgroup and every prefix is a C-marking.

Induct on length. A marking of length $m+1$ has a prefix equivalent
under $A$ to one already visited. After making that change, its
last entry is a candidate passing the exact extension test. An
element of the prefix stabilizer sends it to the chosen orbit
representative, so its full orbit is visited. Conversely, an
automorphism between two visited extensions first identifies
their prefixes. The induction hypothesis makes those prefixes
equal; the automorphism then belongs to $A_t$, so their last
entries have the same candidate orbit. This proves uniqueness.
Independence forbids any extension once the prefix generates $G$.
The fixed depth and finite candidate sets prove termination.
\end{proof}

\begin{lemma}[A necessary completion test]\label{completion}
Every extension of $t$ generates a subgroup of
$K(t)=\gen{H(t),P(t)}$. If a property required of a recorded
marking cannot hold in any subgroup of $K(t)$, its whole branch
may be discarded. In particular, when $G$ acts on $\Omega$,
an intransitive $K(t)$ has no transitive completion. A proper
$K(t)<G$ has no completion generating $G$.
\end{lemma}
\begin{proof}
Every future entry has index at least $m+1$, hence commutes with
the first $m-1$ old marks and belongs to $P(t)$. All completed
generators therefore lie in $K(t)$. A subgroup of an intransitive
group is intransitive, and a subgroup of a proper group is proper.
The conclusion follows. After appending $x$, the new pool is
$P(t)\cap C_G(t_m)$, independently of $x$. Caching these complete
commuting sets therefore preserves the recursion exactly.
The subgroup test must use the entire pool, not merely orbit
representatives. Pruning changes the counts of visited prefixes;
these are not then a census of all partial C-markings.
\end{proof}

\paragraph{Permutation application.}
Take $G=A_n$ and $A=S_n$, acting by simultaneous conjugation.
This action classifies the actual permutation markings up to
renumbering the points. Its stabilizer at $t$ is the successive
centralizer $C_{S_n}(t_1,\ldots,t_m)$. All even involutions are
included: in the implementation they are the disjoint union of
all order-two conjugacy classes of $A_n$. Every transitive
C-marking under consideration consequently occurs, with no
classification of transitive subgroups assumed. If a visited
proper subgroup is already transitive, it is recorded and still
extended. Transitivity is not a stopping condition.

Permutation equivalence differs from abstract marked-group
isomorphism. The latter is tested by the homomorphism sending
the ordered generators of one generated group to those of the
other, with an explicit bijectivity test. Reversing the tuple
is duality. An abstractly self-dual marking can have a permutation
representation inequivalent to its reverse.

\paragraph{Attribution and computational premises.}
Enumeration by prefix stabilizers is established methodology:
Leemans--Mixer, \emph{Algorithms for classifying regular polytopes
with a fixed automorphism group}, Contributions to Discrete
Mathematics 7(2) (2012), Section 5, attributes it to Hartley--Hulpke
(2010). The present proof specifies an ambient subgroup search,
the recursive intersection test and the necessary completion
test used by the accompanying program. No priority for the
general orbit method is asserted.

The complete GAP program and exact execution are required finite
premises for each application, including its ambient model,
all involutions, candidate orbits and final tuples. Exact group
generation, conjugacy classes, centralizers, subgroup intersection,
orbits and homomorphism checks retain GAP's software correctness
as an explicit premise. The polytope correspondence and the
end-intersection criterion remain mathematical inputs. A saved
count without a complete execution is not a census certificate;
no Lean verification is claimed here.
\end{document}
